% 論文用: scenario / helper / rpl / header / tables の関係
% \usepackage{tikz}\usetikzlibrary{arrows.meta,calc,positioning}
\begin{figure}[htbp]
\centering
\begin{tikzpicture}[
  font=\small,
  >=Stealth,
  box/.style={
    draw,
    rounded corners=1.2pt,
    align=center,
    inner sep=5pt,
    minimum height=1.2cm,
    text width=3.6cm
  },
  wide/.style={box, text width=7.6cm},
  arr/.style={->, thick}
]
  \node[wide] (scenario) {\textbf{scenario}\\[2pt]
    {\scriptsize\texttt{example-ping-...-rpl.cc}}\\
    {\scriptsize 実験の組み立て}};

  \node[wide, below=0.5cm of scenario] (helper) {\textbf{helper}\\[2pt]
    {\scriptsize\texttt{rpl-helper.h/.cc}}\\
    {\scriptsize ノードへプロトコルを載せる}};

  \node[wide, below=0.5cm of helper] (rpl) {\textbf{rpl}\\[2pt]
    {\scriptsize\texttt{rpl.h/.cc}}\\
    {\scriptsize 動作の本体（送受信・親選択・転送）}};

  \node[box, below left=0.65cm and 0.05cm of rpl] (header) {\textbf{header}\\[2pt]
    {\scriptsize\texttt{rpl-header.h/.cc}}\\
    {\scriptsize メッセージの型}};

  \node[box, below right=0.65cm and 0.05cm of rpl] (tables) {\textbf{tables}\\[2pt]
    {\scriptsize\texttt{rpl-tables.h/.cc}}\\
    {\scriptsize 近隣と経路の記憶}};

  \draw[arr] (scenario) -- node[right, font=\scriptsize] {使う} (helper);
  \draw[arr] (helper) -- node[right, font=\scriptsize] {生成する} (rpl);
  \draw[arr] (rpl.south) -- ++(0,-0.16) -| node[pos=0.25, font=\scriptsize, left] {載せる} (header.north);
  \draw[arr] (rpl.south) -- ++(0,-0.16) -| node[pos=0.25, font=\scriptsize, right] {読み書き} (tables.north);
\end{tikzpicture}
\caption{scenario, helper, rpl, header, tables の関係。}
\label{fig:rpl-roles}
\end{figure}
